\section{Conclusion}
In many lattice gauge theory applications hadrons carrying high
momenta are required. Due to the exponential increase of relative
errors of $n$-point functions with Euclidean time distances and diminishing
ground state sampling, high momenta previously were very difficult
or impossible to achieve. In Sec.~\ref{sec:basic}
we have introduced a new class of quark smearing methods
for the construction of hadronic interpolators that address and substantially
mitigate these problems. One particular realization of these
methods, that is trivial to implement and comes with very little computational
overhead, is momentum Wuppertal smearing, defined in Eq.~\eqref{fermsmear2}.
We tested this very successfully in Sec.~\ref{sec:momtest}, enabling us to
determine pion and nucleon energies for momenta as high as almost $2\units{GeV}$
and $3\units{GeV}$, respectively, with just 200 measurements, see
Figs.~\ref{fig:disppi} and \ref{fig:dispN}. These figures
also include a comparison
with the conventional method.

In Sec.~\ref{sec:boo} we investigated
the possibility of introducing an (additional)
Lorentz boost~\cite{Roberts:2012tp,DellaMorte:2012xc}. 
With and without a momentum phase factor included,
this gave some improvement over unboosted Wuppertal smearing,
possibly due to a dampening of the phase mismatch by
the more rapid fall-off of the interpolating wave function
in the direction of the momentum.
However, the results obtained were inferior to those of
momentum Wuppertal smearing, without any boost applied.
We conclude that the intuition of contracting the
wave function may be unjustified
since the hadron is moving in real time but not in imaginary
(Euclidean) time.

Iterative methods
(momentum smearing or not) suffer from high iteration
counts $n\propto a^{-2}$ as the continuum limit $a\rightarrow 0$
is approached, in addition to the naive volume factor due to an increasing
number of lattice sites. Therefore, in Sec.~\ref{sec:improve}
we introduce other non-iterative (momentum) smearing methods
that may be more suitable for small lattice spacings
and that also allow for the construction of non-Gaussian shapes.
These will be used by us in the near future.

Realizing momenta that are much larger than the hadron masses of
interest is of fundamental importance in several modern applications,
e.g., in direct determinations of (quasi) distribution
amplitudes~\cite{Aglietti:1998ur,Abada:2001if,Braun:2007wv},
of (quasi) (generalized) parton
distributions~\cite{Ji:2013dva,Ji:2015jwa,Lin:2014zya,Alexandrou:2015rja}
and moments of transverse momentum
distributions~\cite{Hagler:2009mb,Musch:2010ka,Musch:2011er,Engelhardt:2015xja}.
The new method allowed us to extract nucleon masses (employing very
moderate computational resources) up to momenta $\vect{p}^2\approx 7.9\units{GeV}^2$.
This clearly makes the above observables amenable to lattice simulations
in a realistic setting. For three-point functions, these methods
potentially even allow for virtualities
$Q^2=4\vect{p}^2\approx 30\units{GeV}^2$, switching a source momentum of
$-\vect{p}$ into a sink momentum of $+\vect{p}$.
The computation of these quantities with the new smearing
is, depending on the observable, either in progress or planned.\\
